\section{引言}
\subsection{问题背景}

以大语言模型（Large Language Models, LLMs）为代表的人工智能技术在近年来取得飞速发展，其能力的跃升高度依赖于模型规模、训练数据规模、数据质量与领域配比等多维资源的协同配置。为直观理解这一过程，可将大模型训练类比为培养一名高质量学生：参数量 $N$ 相当于脑容量，训练数据量 $D$ 相当于阅读量，计算量 $C$ 相当于教育经费，数据质量 $Q$ 相当于教材质量，领域配比 $\bm{p}$ 相当于各科课时分配。标度律（Scaling Laws）刻画了模型损失随规模变化的经验关系，是理解大模型能力形成的核心工具\cite{kaplan2020scaling,hoffmann2022chinchilla}。

然而，随着人工智能对高质量语料的需求持续增长，干净数据日益稀缺、数据污染问题日趋严重，数据墙（Data Wall）问题日益突出\cite{villalobos2022data}。2026 年《Nature》发表的研究表明，一个仅有 80 亿参数的小规模模型 OpenScholar-8B 在引文准确率、降低幻觉等任务上的表现优于 GPT-4o\cite{asai2026openscholar,brown2020gpt3}，充分说明数据质量与领域配比对模型能力具有不可忽视的作用。扩大模型规模、提升数据质量与延长上下文长度均需付出算力代价，而算力始终有限，四者之间存在复杂的取舍关系。

因此，如何利用真实公开数据建立数学模型，刻画参数规模、数据规模、数据质量、领域配比与交叉熵损失之间的定量关系，并在算力约束下求解最优资源配置，成为亟待解决的关键科学问题。本赛题要求参赛队伍在无须底层系统开发与大规模训练经验的前提下，完成这一建模任务，为算力受限条件下的大模型能力提升提供理论依据与决策参考。

\subsection{问题重述}

\textbf{问题1：} 对 22 个质量指标统一方向并建立综合评价模型，给出样本级、语料级与领域级质量评分 $Q$，定义并消解质量指标间的冲突，并建立 17 域配比与交叉熵损失之间的定量关系。

\textbf{问题2：} 将数据质量 $Q$ 与领域配比 $\bm{p}$ 纳入经典标度律，建立同时包含 $N$、$D$、$Q$、$\bm{p}$ 的广义标度律，分析各因素边际效用与弹性，推导质量与参数规模可相互替代的条件。

\textbf{问题3：} 在算力预算约束下，联合优化参数量 $N$、数据量 $D$、领域配比 $\bm{p}$ 与数据质量 $Q$，并分析预算跨越不同量级时资源配置是否发生结构性转移。

\textbf{问题4：} 结合历史评测数据建立技术演进模型，分离并量化规模扩张与非规模技术进步的贡献占比，预测未来 12 至 24 个月开源大语言模型能力的前沿边界。
